\begin{figure}[!htb]
\centering
\begin{tikzpicture}[x=1cm,y=1cm,font=\small,>=Latex]
 \node[draw,rounded corners,fill=gray!8,minimum width=4.3cm] (prefix) at (0,3.15)
   {common choices at strict ancestors};
 \node[draw,circle,inner sep=3pt] (u) at (0,1.85) {$U$};
 \draw[->,thick] (prefix)--node[right] {stable earlier keys}(u);
 \node[draw,rounded corners,fill=blue!8] (childi) at (-3,.15) {$U_i$, $i<M$};
 \node[draw,rounded corners,fill=orange!12] (childm) at (3,.15) {$U_M$};
 \draw[->,thick,blue!70!black] (u)--node[above left,align=right]
   {$y_n(t)$: selector $i$ succeeds\\earlier children rejected}(childi);
 \draw[->,thick,orange!80!black] (u)--node[above right,align=left]
   {$x$: all $M-1$ selectors vanish\\default choice}(childm);
 \node[align=center,below=5mm of childi] (terminal)
   {local route to $L_i(y_n(t))$\\terminal entry equals $1$};
 \draw[->,blue!70!black] (childi)--(terminal);
\end{tikzpicture}
\caption{Route preservation for $x\in G$, $n\in W_{(U,i)}$ and
$t\in[1,2]$, conditional on $Q_i(y_n(t))=1$.
The center $x$ and its translate
$y_n(t)$ make the same required choices above the first default node $U$.
At $U$, the center takes child $M$, while the translate rejects the earlier
children and takes child $i$. The local terminal success then places
$y_n(t)$ in $B$. The diagram is conditional on the successful local
predicate; it does not assume independence of the chosen leaves.}
\label{fig:routing}
\end{figure}
